% Generated draft: requires mathematical and source review.
\documentclass[11pt]{article}
\usepackage[margin=25mm]{geometry}
\usepackage{fontspec}
\setmainfont{Latin Modern Roman}
\usepackage{amsmath,amssymb,mathtools,unicode-math}
\setmathfont{Latin Modern Math}
\usepackage{enumitem,needspace,xurl,hyperref}
\setlength{\parindent}{0pt}
\setlength{\parskip}{5pt}
\setlength{\emergencystretch}{3em}
\newcommand{\N}{\mathbb{N}}
\newcommand{\Z}{\mathbb{Z}}
\newcommand{\Q}{\mathbb{Q}}
\newcommand{\R}{\mathbb{R}}
\newcommand{\C}{\mathbb{C}}
\newcommand{\F}{\mathbb{F}}
\newcommand{\A}{\mathbb{A}}
\newcommand{\PP}{\mathbb{P}}
\newcommand{\E}{\mathbb{E}}
\newcommand{\eps}{\varepsilon}
\newcommand{\dd}{\,\mathrm{d}}
\begin{document}
\section*{UnsolvedMath ID: 30003544}
\textbf{Source problem code:} catalogue entry
\subsection{Definitions and mathematical statement}
Let \(M\) be a smooth closed manifold of real dimension \(4\), where closed means compact and without boundary. An almost Kähler structure \((g,J,\omega)\) consists of a Riemannian metric \(g\), a smooth endomorphism \(J:TM\to TM\) satisfying \(J^2=-\mathrm{id}_{TM}\) and \(g(JX,JY)=g(X,Y)\), and the two-form \(\omega(X,Y)=g(JX,Y)\), which satisfies \(\mathrm{d}\omega=0\). These identities hold for all smooth vector fields \(X,Y\). The metric is Einstein if \(\operatorname{Ric}_g=\lambda g\) for a real constant \(\lambda\). The almost complex structure \(J\) is integrable if it arises from complex coordinate charts; equivalently, its Nijenhuis tensor vanishes:
\[
N_J(X,Y)=[JX,JY]-J[JX,Y]-J[X,JY]-[X,Y]=0
\]
for all smooth vector fields \(X,Y\).

\textbf{Question (Goldberg conjecture in four dimensions).} For every closed almost Kähler four-manifold \((M,g,J,\omega)\), does
\[
\operatorname{Ric}_g=\lambda g\quad\Longrightarrow\quad N_J\equiv 0
\]
hold? There is no restriction on the sign of \(\lambda\), and no additional curvature hypothesis is imposed. Equivalently, must every such Einstein almost Kähler structure be Kähler?

The supplied dated curation reports a known nonnegative-scalar-curvature theorem attributed to Sekigawa and several results under additional curvature assumptions, while reporting the general compact Einstein almost Kähler Goldberg conjecture, including the negative-scalar-curvature case, as open. These are supplied literature qualifications, not independently verified assertions here; the target statement above is specifically four-dimensional.
\subsection{Short English statement}
Must the almost complex structure of every compact, boundaryless, real four-dimensional Einstein almost Kähler manifold be integrable? The question allows scalar curvature of either sign or zero. The supplied record labels the problem open; this draft does not independently verify that status.
\subsection{Sources}
\begin{itemize}
\item Local UnsolvedMath registry, collection \emph{Oberwolfach Reports Open Problems}, permanent ID \(30003544\), exact source code \texttt{catalogue entry}, title \emph{Goldberg Conjecture in Four Dimensions}. The record gives proposed year 2017 and supplies no named proposer. Its statement status is \texttt{corrected\_verified}, described as preserving mathematical content while repairing notation, typography, or omitted definitions; this is a supplied registry designation.
\item Supplied source citation: \emph{Analysis, Geometry and Topology of Positive Scalar Curvature Metrics} (2018), \emph{Oberwolfach Reports}, pp.\ 2223--2298, DOI \texttt{10.4171/owr/2017/36}. Original OWR report: \url{https://doi.org/10.4171/owr/2017/36}. The background describes these reports as documenting mathematical workshops and their open-problem discussions. The citation year 2018 and the record's proposed year 2017 are retained without resolving their relationship.
\item Supplied literature source: Conti--Fernandez, described in the record as an almost-cokähler analogue: \url{https://arxiv.org/abs/1409.4437}. This is supplied as related literature, not as a resolution of the stated four-dimensional problem.
\item Literature assessment date: \texttt{2026-08-21}. The supplied dated curation labels the problem open, reports Sekigawa's nonnegative-scalar-curvature theorem and several results under additional curvature assumptions, and describes the full compact Einstein almost Kähler Goldberg conjecture, including negative scalar curvature, as remaining open. The record expressly characterizes this as dated literature triage derived from the supplied OWR curation and calls for independent verification of time-sensitive status and cited claims. This assessment has not been independently verified here.
\end{itemize}
\end{document}
